%%=============================================================================
%% Inleiding
%%=============================================================================

\chapter{\IfLanguageName{dutch}{Inleiding}{Introduction}}%
\label{ch:inleiding}

% De inleiding moet de lezer net genoeg informatie verschaffen om het onderwerp te begrijpen en in te zien waarom de onderzoeksvraag de moeite waard is om te onderzoeken. In de inleiding ga je literatuurverwijzingen beperken, zodat de tekst vlot leesbaar blijft. Je kan de inleiding verder onderverdelen in secties als dit de tekst verduidelijkt. Zaken die aan bod kunnen komen in de inleiding~\autocite{Pollefliet2011}:

% \begin{itemize}
%   \item context, achtergrond
%   \item afbakenen van het onderwerp
%   \item verantwoording van het onderwerp, methodologie
%   \item probleemstelling
%   \item onderzoeksdoelstelling
%   \item onderzoeksvraag
%   \item \ldots
% \end{itemize}

\section{Context, achtergrond}%
\label{sec:inl_context}

Binnen de opleiding \emph{Toegepaste Informatica} op \emph{HOGENT} hebben de meeste \emph{OLODs} (opleidingsonderdelen) gelinkt aan het keuzeprofiel ``\emph{IT infrastructure engineer}'' de opzet om het beheer van (virtuele) serveromgevingen aan te leren. Bij het opzetten van dit soort infrastructuur is helderheid over elke stap van het proces en communicatie met elke (al dan niet technisch ingestelde) stakeholder cruciaal.
Om deze duidelijkheid te waarborgen, wordt er binnen de opleiding (meerbepaald tijdens de analyse-\emph{OLODs}) ook aangeleerd om te werken met \emph{UML}-modeling software zoals \emph{Visual Paradigm}\footnote{\url{https://www.visual-paradigm.com/}} en \emph{draw.io}\footnote{\url{https://www.drawio.com/}}. Deze software wordt gebruikt om diagrammen te ontwerpen zoals Activity Diagrams, System Sequence Diagrams, domeinmodellen en BPMN-diagrammen. Deze diagrammen zijn hoofdzakelijk gericht op softwareontwikkeling.

Met deze software kan men ook netwerk- en implementatiediagrammen ontwerpen, die de specificaties, software-elementen (databanken, webservers, firewalls, \ldots) en verbanden tussen verschillende servers in een bepaalde infrastructuur verduidelijken. Deze servers worden op basis van deze diagrammen daarna respectievelijk opgezet via \emph{Infrastructure-as-Code}-tools (zoals \emph{Terraform}\footnote{\url{https://developer.hashicorp.com/terraform}}, \emph{OpenTofu}\footnote{\url{https://opentofu.org/}}, \emph{Vagrant}\footnote{\url{https://developer.hashicorp.com/vagrant}}, \ldots) en geconfigureerd met ``\emph{configuration management}''-tools (\emph{Ansible}\footnote{\url{https://docs.ansible.com/}}, \emph{Puppet}\footnote{\url{https://www.puppet.com/}}, \emph{Chef}\footnote{\url{https://www.chef.io/}}, \ldots). De \emph{HOGENT}-\emph{OLODs} ``\emph{DevOps Project: Operations}'', ``\emph{Cybersecurity Advanced}'' en meerbepaald ``\emph{Infrastructure Automation}'' introduceren ook het gebruik van \emph{Ansible}, één van de leidende ``\emph{configuration management}''-tools.

\section{\IfLanguageName{dutch}{Probleemstelling}{Problem Statement}}%
\label{sec:inl_probleemstelling}

% Uit je probleemstelling moet duidelijk zijn dat je onderzoek een meerwaarde heeft voor een concrete doelgroep. De doelgroep moet goed gedefinieerd en afgelijnd zijn. Doelgroepen als ``bedrijven,'' ``KMO's'', systeembeheerders, enz.~zijn nog te vaag. Als je een lijstje kan maken van de personen/organisaties die een meerwaarde zullen vinden in deze bachelorproef (dit is eigenlijk je steekproefkader), dan is dat een indicatie dat de doelgroep goed gedefinieerd is. Dit kan een enkel bedrijf zijn of zelfs één persoon (je co-promotor/opdrachtgever).

Het huidige proces waarbij (complexe) netwerk- en implementatiediagrammen worden geïnterpreteerd en manueel worden omgezet tot een ``\emph{configuration management}''-ondersteunde serveromgeving kan, volgens \textcite{Nicacio2020}, tijdrovend en vatbaar voor ambiguïteit en foute interpretaties zijn.

Bij implementatiediagrammen kan er onder meer onduidelijkheid onstaan bij (de versies van) software-packages, hun afhankelijkheden/dependencies, de firewall-regels voor elke server en hoe dit het bereik buiten het lokale netwerk beïnvloedt.
Voor netwerkdiagrammen kan dit bijvoorbeeld gaan over hardwarespecificaties van de servers, (de versies van) hun besturingssystemen en de declaratie van de \emph{DNS} server(s).

\section{\IfLanguageName{dutch}{Hoofdonderzoeksvraag}{Research question}}%
\label{sec:inl_hoofdonderzoeksvraag}

% Wees zo concreet mogelijk bij het formuleren van je onderzoeksvraag. Een onderzoeksvraag is trouwens iets waar nog niemand op dit moment een antwoord heeft (voor zover je kan nagaan). Het opzoeken van bestaande informatie (bv. ``welke tools bestaan er voor deze toepassing?'') is dus geen onderzoeksvraag. Je kan de onderzoeksvraag verder specifiëren in deelvragen. Bv.~als je onderzoek gaat over performantiemetingen, dan 

De hoofdonderzoeksvraag voor deze bachelorproef luidt als volgt:

Hoe kan men een gegeven netwerk- en implementatiediagram (in tekstuele vorm via bijvoorbeeld \emph{PlantUML}\footnote{\url{https://plantuml.com/}}) van een infrastructuur met minimale menselijke interventie omzetten tot een overeenkomende, gebruiksklare ``\emph{configuration management}''-ondersteunde (bijvoorbeeld met \emph{Ansible}) omgeving?

Voor deze omgeving wordt ervan uit gegaan dat de eenheden (bare-metal servers, virtuele machines, containers, \ldots) alvorens uitgerold zijn via een \emph{Infrastructure-as-Code}-tool en dat ze enkel nog achteraf geconfigureerd moeten worden via een ``\emph{configuration management}''-tool.

% TODO: Eventueel andere deelvragen vinden of bestaande vervangen, afhankelijk van de inhoud van de komende hoofdstukken

\section{Deelvragen over het probleemdomein}%
\label{sec:inl_deelvragen_probleemdomein}

\begin{itemize}
  \item Wat zijn de doeleinden (business- of technisch-gericht) van \emph{UML}?
  \item Welke infrastructuur-gerelateerde informatie kan een netwerk- en implementatiediagram \emph{niet} weergeven (dat een \emph{Infrastructure-as-Code}- en/of ``\emph{configuration management}''-tool wel kan)?
  \item Welke stappen zijn er tussen het ontwerp van \emph{UML}-diagrammen en de eerste opstelling van een ``\emph{configuration management}''-ondersteunde omgeving?
\end{itemize}

\section{Deelvragen over het oplossingsdomein}%
\label{sec:inl_deelvragen_oplossingsdomein}

\begin{itemize}
  \item Naast \emph{PlantUML} en \emph{MermaidJS}, met welke andere talen kan men \emph{UML}-diagrammen ontwerpen?
  \item Buiten \emph{Ansible}, \emph{Puppet} en \emph{Chef}, wat zijn andere ``\emph{configuration management}''-tools die voor dit soort doeleinden geschikt zijn?
  \item Aangezien er al tools bestaan om \emph{UML}-diagrammen te genereren op basis van ``\emph{configuration management}''-ondersteunde omgevingen, zoals één van \textcite{teramako2023}, welke meerwaarde zou deze oplossing (die in de omgekeerde richting werkt) kunnen bieden?
  \item Is het mogelijk om deze oplossing te ontwikkelen aan de hand van ``reverse engineering'' van de bovenstaande tools?
\end{itemize}

\section{\IfLanguageName{dutch}{Onderzoeksdoelstelling}{Research objective}}%
\label{sec:inl_onderzoeksdoelstelling}

% Wat is het beoogde resultaat van je bachelorproef? Wat zijn de criteria voor succes? Beschrijf die zo concreet mogelijk. Gaat het bv.\ om een proof-of-concept, een prototype, een verslag met aanbevelingen, een vergelijkende studie, enz.

De doelstelling van deze bachelorproef is om een proof-of-concept te ontwikkelen voor een converter (of omzetter) van \emph{PlantUML}-diagrammen naar \emph{Ansible}-code. \emph{PlantUML} en \emph{Ansible} zijn de twee tools die op dit moment voor ogen staan, maar er zullen tijdens het onderzoek ook alternatieven gezocht en geëvalueerd worden.

Deze oplossing zal een meerwaarde bieden voor educatieve doeleinden binnen de opleiding \emph{Toegepaste Informatica}. Neem als voorbeeld een \emph{OLOD} zoals ``\emph{Cybersecurity Advanced}'' -- een vak waarin studenten leren om de beveiliging van netwerken te versterken aan de hand van virtuele labo's. De tech-industrie evolueert snel en de labo's die voor dit \emph{OLOD} ontworpen zijn, zullen daardoor ook snel verouderen.

Aangezien met deze oplossing de drempel tussen de ontwerp- en implementatiefase verkleint, zouden docenten sneller nieuwe (virtuele) labo-omgevingen kunnen ontwerpen, met (in theorie) onmiddelijk beschikbare ``\emph{configuration management}''-code om deze meteen te configureren. Voor studenten kan dit ook een kans bieden om potentiële infrastructuur-gerelateerde experimenten tijdens hun studies sneller effectief vorm te geven.

\section{\IfLanguageName{dutch}{Opzet van deze bachelorproef}{Structure of this bachelor thesis}}%
\label{sec:inl_opzet-bachelorproef}

% Het is gebruikelijk aan het einde van de inleiding een overzicht te
% geven van de opbouw van de rest van de tekst. Deze sectie bevat al een aanzet
% die je kan aanvullen/aanpassen in functie van je eigen tekst.

De rest van deze bachelorproef is als volgt opgebouwd:

In Hoofdstuk~\ref{ch:literatuurstudie} (\nameref{ch:literatuurstudie}) wordt een overzicht gegeven van de stand van zaken binnen het onderzoeksdomein, op basis van een literatuurstudie.

In Hoofdstuk~\ref{ch:methodologie} (\nameref{ch:methodologie}) worden de verschillende fases van dit onderzoek toegelicht, alsook de gebruikte onderzoekstechnieken om antwoorden te kunnen formuleren op de hoofdonderzoeksvraag en deelvragen.

% Vul hier aan voor je eigen hoofstukken, één of twee zinnen per hoofdstuk

In Hoofdstuk~\ref{ch:manuele_testomgeving} (\nameref{ch:manuele_testomgeving}) wordt het ontwerp en de implementatie van de manuele testomgeving voor de proof-of-concept beschreven. Dit bevat een evaluatie van bestaande tools in de betrokken vakgebieden (\emph{UML} en \emph{configuration management}), alsook het ontwerp van de verwachte invoer en uitvoer voor de converter.

In Hoofdstuk~\ref{ch:poc} (\nameref{ch:poc}) wordt toegelicht hoe de proof-of-concept, in de vorm van een converter voor \emph{PlantUML} naar \emph{Ansible} (\emph{PlantUML2Ansible}), is ontwikkeld.

In Hoofdstuk~\ref{ch:testcases} (\nameref{ch:testcases}) wordt de definitie van de verschillende testcases, waar de converter op wordt getoetst, beschreven. De resultaten van deze testen worden ook in dit hoofdstuk toegelicht.

In Hoofdstuk~\ref{ch:conclusie_discussie} (\nameref{ch:conclusie_discussie}), ten slotte, wordt de conclusie gegeven en een antwoord geformuleerd op de onderzoeksvragen. Daarnaast wordt in de discussie ook een aanzet gegeven voor toekomstig onderzoek binnen dit domein.